\makeabstract
{ % #1: Title
Multitargeting CYP3A4, PI3K\ensuremath{\alpha}, and CDK4/6 Overcomes Hormone Therapy Resistance in Mutant PI3K\ensuremath{\alpha} ER+/HER2\ensuremath{-} Breast Cancer
}
{ % #2: Authorship
Elizabeth Koa\textsuperscript{1, 2},
Zhijun Guo\textsuperscript{1}, 
Jianxun Lei\textsuperscript{1}, 
Antonio D'Assoro\textsuperscript{3}, 
Matthew Goetz\textsuperscript{3}, 
David Potter\textsuperscript{1}
}
{ % #3: Affiliations (use \\ to start a newline)
\textsuperscript{1} Division of Hematology, Oncology, and Transplantation, University of Minnesota, Minneapolis, MN\\ 
\textsuperscript{2} Amherst College, Amherst, MA\\ 
\textsuperscript{3} Mayo Clinic, Rochester, MN
}
{ % #4: Purpose
About 40\% of patients with estrogen receptor-positive HER2 negative (ER+ HER2-) metastatic breast cancer (MBC) exhibit hormone therapy resistance (HTR) driven by mutant PI3K\ensuremath{\alpha}, limiting the effectiveness of endocrine therapies such as the selective estrogen receptor degrader (SERD) fulvestrant (F). Mutant PI3K HTR is associated with an increase in the arachidonic acid (AA) precursor of cancer promoting -6 epoxy polyunsaturated fatty acids (-6 EpFAs) and upregulation of the enzyme CYP3A4 that synthesizes EpFAs from AA. In the first part of this study, we measured levels of various proteins, including CYP3A4, involved in breast cancer signaling pathways, to find which proteins are upregulated in HTR cells. In the second part of this study, we tested different drug combinations in a synergy study to find out how to overcome hormone therapy resistance. 
}
{ % #5: Methods
F-selected MCF-7AC1 clonal cell lines were screened and the rate of CYP3A4 upregulation was measured using a western blot. Next, we treated model HTR cells with the mutant PI3K\ensuremath{\alpha} inhibitor tersolisib (T) and the CYP3A4 inhibitor C5F2-HCB and measured Akt, ERK, and S6 kinase phosphorylation. We also treated model HTR cells with the cell cycle inhibitor abemaciclib (A) and C5F2-HCB. Here, we measured Cyclin E1 and CDK6 levels. Lastly, we tested the combinations A+T, A+T+F, and A+T+F+C5F2-HCB in a MTT viability assay. The computer program CompuSyn produced combination index values (CI values) at different effective doses.
}
{ % #6: Results
CYP3A4 was upregulated in 22.2\% of cells. Treatment with T significantly decreased Akt phosphorylation. Treatment with C5F2-HCB increased Akt and decreased ERK and S6 Kinase phosphorylation. Surprisingly, we found that treatment with A increased Cyclin E1 Levels and treatment with C5F2-HCB decreased CDK6 levels. C5F2-HCB synergized with a combination of the mutant PI3K\ensuremath{\alpha} inhibitor T, the CDKi A, and F (CI = 0.27). When A, T, F, and C5F2-HCB are combined, there was higher cell inhibition for all drug concentrations compared to when C5F2-HCB was administered alone. The quadruple combination showed the best synergy, with all CI values significantly less than one. 
}
{ % #7: Conclusion
C5F2-HCB, a compound our lab has made previously that inhibits CYP3A4, could be used in multitargeted drug therapies. Treatment with C5F2-HCB decreased CDK6 levels suggesting that C5F2-HCB may overcome resistance to CDKis such as P and A in HTR mutant PI3K\ensuremath{\alpha} cells. Finally, a combination of 4 drugs---C5F2-HCB (CYP3A4i), abemaciclib (CDKi), fulvestrant, and tersolisib (mutant PI3K\ensuremath{\alpha}i)---show strong synergy. This means multitargeted inhibition of CYP3A4, CDKs, and mutant PI3K\ensuremath{\alpha} not only overcomes HTR, but suggests decreased toxicity and increased treatment efficacy for patients in the clinic. This promising combination supports its further evaluation in preclinical models.
}
 \newpage